\section*{Abstract}
\addcontentsline{toc}{section}{Abstract}

Yaound\'e is a multilingual city in which ordinary conversation mixes French, English,
Pidgin, Fulfulde and Ewondo with locally coined slang. This report documents
\textbf{Francanglais Studio}, a complete compiler front end --- scanner, lexer, grammar
preparation, LL(1) table construction and a table-driven predictive parser --- built to
perform lexical and syntactic analysis on speech actually collected in the city, and
wrapped in a web application so that every compiler artefact is directly inspectable.

The team collected \CorpusTotal\ statements by hand across taxis, markets, roadside
businesses, neighbourhoods and the university, cross-checked each transcription, and
recorded the attestation as a software-enforced invariant rather than a claim in prose.
From that data we built a \LexiconTotal-entry lexicon spanning \OriginCount\ donor
languages and \TerminalCount\ terminal classes, authored a \ProductionCount-production
context-free grammar in its natural left-recursive form, and mechanically transformed it
in \LedgerSteps\ recorded steps into an LL(1) grammar whose parsing table is verified
conflict-free by an automated test. The resulting parser accepts \CorpusAccepted\ of the
\CorpusTotal\ collected statements; the \CorpusRejected\ rejections are documented down to
the specific LL(1) conflict that prevents a fix rather than papered over.

This document is organised as a combined Software Requirements Specification
(Part~\ref{part:srs}) and Software Design Document (Part~\ref{part:sdd}), followed by the
compiler construction work itself (Part~IV) and its evaluation (Part~V). Every table of
data it contains is exported directly from the running analyzer by a build script, so the
document cannot drift from the software it describes.

\vspace{0.8em}
\begin{keypoint}[Reproducibility]
No figure or table in this report was transcribed by hand. Running
\texttt{build.ps1} regenerates the data tables from the live analyzer, re-renders the
PlantUML diagrams, and recompiles this PDF. Appendix~\ref{app:repro} gives the procedure.
\end{keypoint}
